\chapter{Conclusão}
\label{cap:conclusao}

\section{Síntese das Contribuições}

O trabalho propôs e implementou um Data Warehouse para a análise da execução orçamentária de Juiz de Fora, com o objetivo de tornar os dados públicos mais acessíveis e oferecer um modelo replicável para outros municípios.

O primeiro objetivo previa a coleta e documentação de dados de receitas e despesas do Portal da Transparência da Prefeitura. Foi atendido: o \textit{pipeline} de extração consome seis fontes (despesa mensal consolidada, receita mensal comparativa, receita mensal prevista, ementário de natureza de receita da STN, função LOA e subfunção LOA do exercício 2026), com tratamento explícito de heterogeneidades como mudança de \textit{layout} no meio do histórico, cabeçalho de planilha sem posição fixa e OCR nos PDFs da LOA.

O segundo objetivo, que previa a modelagem dimensional conforme o modelo de Kimball, foi cumprido. O \textit{warehouse} conta com dez dimensões e três tabelas de fatos organizadas em esquema de estrela compartilhando \texttt{dim\_tempo}. Cada decisão de granularidade foi escolhida em função do que a fonte de fato publica, sem inventar atomicidade que não está presente nos dados originais.

O terceiro objetivo previa a implementação de processos de ETL e foi atendido com Dagster como orquestrador, Python e \texttt{pandas} para extração, DuckDB como banco de dados analítico e dbt para modelagem em SQL. A unidade de idempotência foi definida como uma partição temporal, e três mecanismos garantem que a reexecução dessa mesma partição produza o mesmo estado final do \textit{warehouse}.

O quarto objetivo, que previa a construção de \textit{dashboards} interativos, também foi alcançado. Sete páginas em Streamlit consomem o \textit{warehouse}, estão publicadas em uma URL pública e respondem perguntas como a execução orçamentária por área, a evolução temporal de receitas, a concentração de fornecedores e o desvio entre previsão e arrecadação. A exploração sistemática descrita no Capítulo~\ref{cap:resultados} destacou duas leituras recorrentes: dependência de transferências intergovernamentais na arrecadação e a distância entre empenho e pagamento ao longo do exercício.

Complementarmente, o Capítulo~\ref{cap:avaliacao} compara o portal e o \textit{dashboard} em seis tarefas de fiscalização. O painel reduz os passos manuais de consolidação em todas elas. A lacuna técnica dos portais foi a mais atendida, e a pedagógica, parcialmente.

\section{Considerações Finais}

O processo de construção deste \textit{warehouse} revelou algo que não estava previsto no plano inicial. O problema central da transparência orçamentária municipal não é técnico no sentido mais estreito da palavra.

O problema está na distância entre o que o portal publica e o que o cidadão precisa para fiscalizar. O portal publica relatórios mensais consolidados em planilhas Excel. O cidadão precisa de séries temporais comparáveis, classificações inteligíveis e indicadores que dialoguem com sua experiência cotidiana. Reduzir essa distância é trabalho de arquitetura de dados, e arquitetura de dados é, em última análise, uma decisão sobre quem deve interpretar a informação pública. Quando se publica o valor apenas em colunas mensais, sem a data de cada pagamento e sem abrir o empenho por itens, decide-se o que o cidadão pode ver. Ele enxerga quanto saiu no mês e a descrição do empenho. Não enxerga o ritmo dos pagamentos, nem a quantidade nem o preço do que foi comprado.

A leitura freireana ganha contorno técnico nesse ponto. Conscientização orçamentária, traduzida para o trabalho de quem constrói \textit{data warehouses}, significa modelar o dado pensando em quem vai lê-lo, e não apenas em quem vai gerá-lo.

\section{Trabalhos Futuros}

A continuação natural do trabalho segue cinco direções concretas:

(1) Integração da dotação autorizada pela LOA, viabilizando indicadores percentuais de execução sobre o autorizado, que são a métrica padrão utilizada por órgãos de controle. A classificação funcional do exercício de 2026 já está no DW.

(2) Desagregação da funcional programática em programa e ação, com \textit{parsers} LOA adaptados a exercícios anteriores e, se necessário, aos ementários oficiais da STN para níveis ainda ausentes.

(3) Captura de licitações e contratos como processos de negócio adicionais, permitindo rastrear a cadeia completa do gasto e analisar o tempo médio entre as etapas.

(4) Implementação de SCDs do tipo 2 nas dimensões sujeitas a mudanças históricas, permitindo análises "como era" \textit{vs.} "como é".

(5) Replicação do modelo em pelo menos um município de porte semelhante, com análise comparativa, para validar a generalidade do esquema dimensional proposto e identificar pontos de adaptação locais.

Cada uma dessas direções é viável com o ferramental atual e não exige reescrita das camadas já implementadas. O DW aqui construído, quando tomado em conjunto com a discussão teórica que o sustenta, não esgota o problema da transparência orçamentária. Oferece, em vez disso, uma peça de evidência de que é possível, no escopo de um trabalho, trilhar o caminho entre o dado bruto que a prefeitura publica e a pergunta que o cidadão de fato faria. Essa ponte é a contribuição.
